\centerline{\bf \Huge Тренировочная олимпиада II}
\centerline{\bf \Huge }

\begin{flushright}
        \scriptsize{}
\end{flushright}

\problem{1}
Сколько существует натуральных чисел $n$, не превосходящих 2022, таких, что $(n+1)(n+ 2)(n+3)$ делится на $343 ?$

\problem{2}
На заводе производят $N$ видов мороженого. На дегустацию пришли 7 детей, каждый из них попробовал 7 видов мороженого. При этом каждый ребёнок попробовал ровно 4 вида, которые не пробовал никто другой, а каждый из оставшихся 3 видов попробовал кто-то ещё. Какое наибольшее значение может принимать $N$, если каждый вид мороженого попробовал хотя бы один ребёнок?

\problem{3}
Алёна вписывает цифры в клетки прямоугольной таблицы, в каждую клетку она вписывает не более одной цифры. Если она впишет 30 цифр, то обязательно найдётся строка, в которой будет записано хотя бы 3 цифры. Если она впишет 15 цифр, то обязательно найдётся столбец, в котором будет записано хотя бы 4 цифры. Какое наибольшее количество клеток может быть в таблице?

\problem{4}
В магазине в наличии есть арбузы и дыни. Сегодня привезли новые арбузы и дыни, а часть арбузов и дынь купили. В результате общее число арбузов и дынь сократилось на $10 \%$, а доля арбузов увеличилась с $55 \%$ до $65 \%$. Какое наименьшее количество арбузов могло остаться?

\problem{5}
Если в компьютер вставить карточку с ненулевым числом $a$, а затем карточку с ненулевым числом $b$, то он вернёт обе эти карточки, а также выдаст новую карточку с числом $1-\frac{a}{b}$. Можно ли, имея карточки с числами 1 и 10, получить карточку с числом 1000?